\chapter{Testing e validazione}
\label{chap:testing-validazione}

% ==================================================
% 6.1 Strategia
% ==================================================

\section{Strategia di verifica}
\label{sec:strategia-testing}


\subsection{Obiettivi della validazione}
\label{subsec:obiettivi-validazione}

% Obiettivo:
% verificare che il prototipo soddisfi i requisiti
% definiti nel Capitolo 3.
%
% La validazione deve dimostrare:
%
% - correttezza statica del progetto;
% - producibilità della build;
% - eseguibilità dei flussi core;
% - correttezza dei comportamenti geografici;
% - applicazione delle misure privacy.
%
% Relazione:
%
% Requisito
% -> Implementazione
% -> Caso di test
% -> Esito


\subsection{Livelli di verifica adottati}
\label{subsec:livelli-testing}

% 1. Verifica statica
%
% - TypeScript
% - ESLint
% - Nx module boundaries
% - Biome
%
% 2. Verifica di build e integrazione
%
% - build Next.js/Nx
% - avvio ambiente
% - database/PostGIS
%
% 3. Validazione funzionale
%
% - scenari utente
% - geolocalizzazione
% - richieste
% - privacy
%
% Dichiarare esplicitamente che il prototipo
% corrente non dispone di una suite automatizzata
% completa di unit/integration/E2E test.


% ==================================================
% 6.2 Qualità statica
% ==================================================

\section{Verifica statica e qualità del codice}
\label{sec:verifica-statica}

\subsection{Type checking}
\label{subsec:typechecking}

% VQ-01 — SUPERATO
%
% Comando:
%
% pnpm exec tsc -p apps/web/tsconfig.json --noEmit
%
% Il controllo si è concluso senza produrre
% errori o diagnostiche TypeScript.


\subsection{Linting e convenzioni}
\label{subsec:linting}

% VQ-02 — SUPERATO
%
% pnpm lint
%
% Nx ha completato correttamente il target lint
% dell'applicazione web.
%
%
% VQ-03 — SUPERATO
%
% pnpm biome:ci
%
% 170 file verificati.
% Nessuna correzione necessaria.


\subsection{Verifica della build}
\label{subsec:build-validation}

% VQ-04 — SUPERATO
%
% pnpm build
%
% Risultati:
%
% - compilazione completata;
% - controllo TypeScript completato;
% - page data raccolti;
% - 14/14 pagine statiche generate;
% - ottimizzazione finale completata;
% - tutte le route applicative riconosciute.


\subsection{Riproducibilità dell'ambiente}
\label{subsec:environment-validation}

% VQ-05 — SUPERATO
%
% pnpm dev:fresh
%
% Verificati con successo:
%
% - installazione dipendenze;
% - avvio PostgreSQL/PostGIS;
% - inizializzazione PostGIS;
% - generazione Prisma Client;
% - sincronizzazione schema;
% - seed database;
% - avvio Next.js.
%
% Osservazione non bloccante:
% configurazione Prisma package.json deprecata
% nelle versioni future.

\section{Validazione dei requisiti funzionali}
\label{sec:validazione-funzionale}


\subsection{Account e autenticazione}
\label{subsec:test-auth}

% TC-AUTH-01
% Registrazione valida.
%
% TC-AUTH-02
% Email duplicata.
%
% TC-AUTH-03
% Conferma email con token valido.
%
% TC-AUTH-04
% Token invalido o scaduto.
%
% Collegare:
%
% RF-01
% RF-02
% RF-03
% RNF-03
% RNF-04


\subsection{Gestione del patrimonio librario}
\label{subsec:test-books}

% TC-BOOK-01
% Creazione libro.
%
% TC-BOOK-02
% Modifica.
%
% TC-BOOK-03
% Eliminazione.
%
% TC-BOOK-04
% Visibilità privata.
%
% TC-BOOK-05
% Gestione ISBN duplicato.


\subsection{Ricerca e discovery}
\label{subsec:test-search}

% TC-SEARCH-01
% Ricerca bibliografica/testuale.
%
% TC-SEARCH-02
% Accesso alla ricerca nearby senza autenticazione.
%
% Atteso:
% redirect al login.


\subsection{Condivisione e richieste}
\label{subsec:test-requests}

% TC-REQ-01
% Creazione richiesta valida.
%
% TC-REQ-02
% Richiesta relativa al proprio libro.
%
% TC-REQ-03
% Richiesta senza autenticazione.
%
% TC-REQ-04
% Libro non disponibile/non valido.
%
% TC-REQ-05
% Aggiornamento stato richiesta.


% ==================================================
% 6.4 Geolocalizzazione
% ==================================================

\section{Validazione della geolocalizzazione e della privacy}
\label{sec:validazione-geolocalizzazione}


\subsection{Geocoding e fallback}
\label{subsec:test-geocoding}

% TC-GEO-01
%
% Geoapify configurato
% -> coordinate.
%
% TC-GEO-02
%
% Geoapify non disponibile
% -> fallback Nominatim.


\subsection{Approssimazione delle coordinate}
\label{subsec:test-coordinate-privacy}

% TC-PRIV-01
%
% Verificare la trasformazione:
%
% coordinate precise
% -> coordinate pubbliche a 3 decimali.
%
%
% TC-PRIV-02
%
% Verificare che discovery e mappa utilizzino
% publicLatitude/publicLongitude
% e non latitude/longitude.


\subsection{Ricerca per prossimità}
\label{subsec:test-nearby}

% Raggi supportati:
%
% 5 km
% 10 km
% 25 km
% 50 km
%
% TC-GEO-03
% Libro interno al raggio.
%
% TC-GEO-04
% Libro esterno al raggio.
%
% TC-GEO-05
% Libro privato escluso.
%
% TC-GEO-06
% Libro indisponibile escluso.
%
% TC-GEO-07
% Ordinamento per distanza.


\subsection{Visualizzazione cartografica}
\label{subsec:test-map}

% TC-MAP-01
% Marker basati su coordinate pubbliche.
%
% TC-MAP-02
% Marker / popup / cluster / modalità 2D-3D.
%
% TC-MAP-03
% Comportamento con prefers-reduced-motion.


% ==================================================
% 6.5 Tracciabilità
% ==================================================

\section{Tracciabilità e sintesi della validazione}
\label{sec:tracciabilita-testing}

% Creare la matrice:
%
% requisito
% -> caso d'uso
% -> caso di test
% -> esito.
%
% Gli esiti devono essere compilati solo
% in seguito all'esecuzione effettiva dei test.

\begin{longtable}{p{0.13\textwidth} p{0.16\textwidth} p{0.27\textwidth} p{0.24\textwidth}}
    \caption{Matrice di tracciabilità della validazione}
    \label{tab:validation-traceability} \\

    \toprule
    \textbf{Requisito} &
    \textbf{Caso d'uso} &
    \textbf{Caso di test} &
    \textbf{Esito} \\
    \midrule
    \endfirsthead

    \toprule
    \textbf{Requisito} &
    \textbf{Caso d'uso} &
    \textbf{Caso di test} &
    \textbf{Esito} \\
    \midrule
    \endhead

    RF-01 & UC-01 & TC-AUTH-01 & Da verificare \\
    RF-02 & UC-01 & TC-AUTH-03 & Da verificare \\
    RF-05 & UC-03 & TC-BOOK-01 & Da verificare \\
    RF-08 & UC-03 & TC-PRIV-01 & Da verificare \\
    RF-12 & UC-06 & TC-GEO-03, TC-GEO-04 & Da verificare \\
    RF-13 & UC-06 & TC-MAP-01, TC-MAP-02 & Da verificare \\
    RF-14 & UC-07 & TC-REQ-01 & Da verificare \\
    RF-15 & UC-07 & TC-REQ-02, TC-REQ-03, TC-REQ-04 & Da verificare \\

    \bottomrule
\end{longtable}